% ===================================================================
% methods_tcad.tex -- Methods: device simulation (Silvaco ATLAS 5.28), measured data, extraction protocol,
% campaign design and pre-registration, scoring/resolution method, execution controls, launch accounting.
% Draft 2026-10-13 (Claude), for human review; no claim here is approved for a manuscript (charter section 9).
% ASCII only. Every number carries a source comment  % <!-- KEY:line -->  (file:line). KEY legend:
%   PKG    = IWO_PHYSICS_CONSTRAINED_MODEL_V1/              PRIV   = iwo-tft-research-private/
%   GEOM   = PKG/config/geometry.yaml                       SOLV   = PKG/config/solver.yaml
%   PHYS   = PKG/config/physics_options.yaml                V3     = PKG/config/iwo_material_model_v3.yaml
%   PROT   = PKG/analysis_2026-10-06_extraction/PROTOCOL.md A1R    = PKG/analysis_2026-10-06_extraction/A1_RESULTS.md
%   PFIG   = PKG/analysis_2026-10-06_extraction/PAPER_FIGURES.md
%   CD     = PKG/analysis_2026-10-10_campaign/CAMPAIGN_DESIGN.md
%   S2PREP = PKG/analysis_2026-10-10_campaign/S2_PREP.md
%   S2FIX  = PKG/analysis_2026-10-10_campaign/S2_SCORER_FIXES_2026-10-12.md
%   S2GATE = PKG/analysis_2026-10-10_campaign/S2_RUNNER_GATES_2026-10-13.md
%   DEVTEX = PKG/report_latex_full/chapters/03_device_data.tex   V2TEX = PKG/report_latex_full/chapters/14_model_v2.tex
%   MAIN   = PKG/report_latex_full/main.tex
%   REG    = PRIV/codex/work_chat/HYPOTHESIS_REGISTER.md    LEDGER = PRIV/03_CLAIMS_LEDGER.yaml
%   DLOG   = PRIV/05_DECISION_LOG.md                        CS     = PRIV/01_CURRENT_STATE.md
% Not opened: PKG/results/runs/run_0075_* and PKG/results/sealed/ (S2-11 is sealed).
% ===================================================================
\subsection{Device simulation}
\label{sec:m-tcad}

\subsubsection{Devices and measured data}
\label{sec:m-tcad-data}

The simulated transistors are the bottom-gate W-doped \InO{} (IWO) thin-film transistors of the thickness series
published by \citet[pp.~2--3]{Januar2026}; the page is the one verified for the stack in
\file{report_latex_full/chapters/03_device_data.tex}. % <!-- DEVTEX:18-19 -->
The devices were made and measured by the fabrication group (National Taiwan University); this work adds no
measurement (decision D-11). % <!-- DLOG:22 ; CS:27 -->
From the gate upwards the stack is a \SI{50}{nm} TiN gate on n$^+$ Si, \SI{15.0}{nm} \HfO{}, \SI{2.0}{nm} \AlO{}, the
sputtered IWO channel (about 2\,\% W, composition basis not stated by the source) and \SI{70}{nm} Pd source and drain
contacts on top of the channel, with no passivation layer. % <!-- GEOM:4-9 ; DEVTEX:23-31 ; V3:19-21 -->
Both oxides were grown by atomic-layer deposition at \SI{250}{\celsius}, and the process ends with a rapid thermal
anneal at \SI{150}{\celsius} for \SI{600}{s}; these enter no parameter. % <!-- DEVTEX:25 ; DEVTEX:33-34 -->
\Cref{tab:m-tcad-stack} lists every structural item with its provenance label.

\begin{table}[htbp]\centering\footnotesize
\caption{Structural items of the simulated device and their provenance labels.}
\label{tab:m-tcad-stack}
\begin{tabularx}{\linewidth}{@{}l>{\raggedright\arraybackslash}Xl@{}}
\toprule
Item & Value & Label \\
\midrule
Gate & TiN \SI{50}{nm} (Si support screened, omitted) & \prov{MEASURED} (schematic) \\ % <!-- GEOM:3-4 -->
Gate dielectric & \HfO{} \SI{15.0}{nm}, $\varepsilon_r = 19.57$ & thickness from schematic; $\varepsilon_r$ \prov{MEASURED} (C--V, SI Fig.~S1) \\ % <!-- GEOM:5 -->
Interlayer & \AlO{} \SI{2.0}{nm}, $\varepsilon_r = 9.0$, touching the channel & thickness from schematic; $\varepsilon_r$ and position \prov{ASSUMED} \\ % <!-- GEOM:6 ; CD:168 -->
Channel & IWO, $t = 2.0$, 6.3, 13.2, \SI{31.8}{nm} (nominal), $\varepsilon_r = 9.3$ & \prov{MEASURED} ($\varepsilon_r$ from C--V, SI Fig.~S1) \\ % <!-- GEOM:7 ; V3:43-46 -->
Source/drain & Pd \SI{70}{nm}, top contacts, ideal Ohmic & geometry from schematic; contact model \prov{ASSUMED} \\ % <!-- GEOM:8 ; V3:209-211 -->
Contact overlap & \SI{2.0}{\micro\metre} per side & \prov{ASSUMED} \\ % <!-- GEOM:14 -->
$L$, $W$ & \SI{20}{\micro\metre}, \SI{290}{\micro\metre} & schematic \\ % <!-- GEOM:11-12 -->
Geometric $\Cox$ & \SI{8.9554e-7}{\farad\per\centi\metre\squared} & \prov{DERIVED} from the rows above \\ % <!-- PROT:20 -->
Gate work function & \SI{4.7}{eV} & \prov{ASSUMED} (TiN 4.5--4.9~eV); degenerate with $\Qf$ and $\chi$ \\ % <!-- V3:206-208 -->
Temperature & \SI{300}{K} & \prov{ASSUMED} (not recorded) \\ % <!-- V3:22-25 -->
Above the channel & air (no cap) & as specified \\ % <!-- GEOM:9 -->
\bottomrule
\end{tabularx}
\end{table}

\paragraph{Measured data.} One transfer curve per film was available: $\Vg$ from $-3$ to $+3$~V in \SI{50}{mV} steps
(121 points) at $\Vd = \SI{0.7}{V}$, for 2.0, 6.3, 13.2 and \SI{31.8}{nm}
(\prov{MEASURED}, claim C001). % <!-- CS:73-76 ; LEDGER:5-9 -->
There is no reverse sweep, no gate current, no second drain bias and only one device per thickness. % <!-- LEDGER:9 ; CS:77 -->
The thicknesses are nominal and the method of their measurement is unknown. % <!-- CS:253 -->
The \SI{31.8}{nm} device is not modelled (decision D-01); it never turns off ($\Id(-3\,\mathrm{V}) =
\SI{3.0e-7}{\Aum}$, claim C029, \prov{DERIVED}) and, having been in the data file since September, it is not a
holdout. % <!-- DLOG:12 ; LEDGER:172-176 -->
Every comparison between simulation and measurement in this work is therefore retrospective. % <!-- REG:37 ; CD:501 -->
The off-state floors, the median current over $\Vg = -2$ to $-0.5$~V, are $4.61\times10^{-15}$, $1.53\times10^{-13}$
and \SI{6.60e-12}{\Aum} for 2.0, 6.3 and \SI{13.2}{nm} (claim C003, \prov{DERIVED}). % <!-- LEDGER:16-20 ; CD:344 -->
\pdPending{P01}{DECISION PENDING: written permission of the fabrication group to use and publish the raw curves
(they are already in the public repository), data custody and authorship; charter section 8}.

\subsubsection{Simulation domain and numerical method}
\label{sec:m-tcad-numerics}

All device simulations use two-dimensional drift--diffusion in Silvaco ATLAS 5.28.1.R driven by DeckBuild 5.0.10.R
\citep{AtlasManual,DeckBuildManual}. % <!-- MAIN:35 -->
The domain contains the physical volumes: air above the channel, the Pd contacts, the \AlO{} and \HfO{} regions and
the TiN gate as a conductor volume. % <!-- GEOM:10 -->
The mesh width is \SI{1}{\micro\metre}, so currents are reported in A/\textmu m. % <!-- GEOM:13 -->
The $x$ mesh is refined to \SI{10}{nm} at the contact edges; along $y$ the spacing is \SI{0.0625}{nm} at the
IWO/\AlO{} interface, $\min(1.0\,\mathrm{nm}, t/8)$ in the film, \SI{0.25}{nm} in \AlO{} and \SI{1.0}{nm} in \HfO{}. % <!-- GEOM:18 ; GEOM:20-23 -->
The interface trap sheet is regularised as a \SI{0.25}{nm} IWO layer; halving that layer changes $\Id$ by
\SI{0.009}{dec}. % <!-- GEOM:19 -->

The solver uses Fermi--Dirac statistics, a Shockley--Read--Hall background and electrons only
(\texttt{carriers=1}); the hole current is numerically zero in every recorded run. % <!-- SOLV:2-3 -->
The Newton scheme includes trap occupancy (\texttt{newton trap}, \texttt{maxtraps=10}, \texttt{itlimit=80},
\texttt{climit=1e-6}, \texttt{ir.tol=1e-19}). % <!-- SOLV:4 -->
In depletion, convergence is accepted with \texttt{XANDRNORM} and a continuity tolerance
$\max(10^{-17}\,\mathrm{A},\ 10^{-5}\times$ the smallest measured active current$)$, up to a measured current of
\SI{1e-9}{\Aum}. % <!-- SOLV:5-8 -->
The trap density of states is integrated on 384/192 levels in the campaign (768/384 in the S1-04 check). % <!-- CD:195 -->
Gate sweeps are stepped \texttt{SOLVE VSTEP/VFINAL} blocks from $-3$~V: \SI{0.2}{V} steps to $-1.0$~V,
\SI{0.05}{V} to \SI{1.5}{V} and \SI{0.1}{V} to \SI{3.0}{V}, after a \SI{0.1}{V} ramp to the sweep start. % <!-- SOLV:10-13 -->
Launches run strictly one at a time (parallel launches failed), with thickness-dependent time limits of 1200, 1500,
2400 and \SI{3000}{s} for 2.0, 6.3, 13.2 and \SI{31.8}{nm} (multiplied per entry where listed). % <!-- CS:285 ; CD:289 ; CD:479 -->
A run passes the current-balance (KCL) check when
$|I_D + I_S + I_G| \le \max(10^{-17}\,\mathrm{A},\ 0.1\times\mathrm{floor}) + 1\,\%\ \sum|I|$. % <!-- CD:356 ; SOLV:19-21 -->
Three tool facts constrain the decks: \texttt{MOBILITY} statements placed between \texttt{SOLVE}s are ignored, at most
ten trap distributions are allowed, and \texttt{DEFECTS} and \texttt{INTDEFECTS} take densities per volume and per
area, respectively. % <!-- CS:288-290 -->

\paragraph{Numerical verification.} Before any hypothesis-scoring launch, the baseline physics was checked under
refinement: mesh $\times0.7$ at \SI{13.2}{nm} (S1-03) changed the active log-current error by $-0.0005$~dec and a
doubled trap-DOS grid at \SI{2.0}{nm} (S1-04) by $-0.00002$~dec, both inside the pre-registered numerics rule
(\SI{0.01}{dec}, \SI{5}{mV} in $\Vthcc$, \SI{3}{mV/dec} in SS). % <!-- REG:29-30 ; CD:298 -->
A native launch (S1-07) confirmed the offline shift/scale transformation used for profiling to a maximum deviation of
$2.05\times10^{-4}$~dec; this validates that transformation only. % <!-- REG:33 -->
The extraction carries a separate numerical budget of $\pm1.5$~mV in $\Vth$, $\pm1$~mV/dec in SS and $\pm0.7\,\%$ in
$\gm$. % <!-- PROT:67 -->

\paragraph{TCAD check rule.} Every change of the deck generator or of a scoring rule is reviewed by Codex and checked
by hand against the ATLAS 5.28 manual, which governs, with the 2004 manual as a second reference; manual pages are
quoted where a model's numerical control is at stake (for example the oxide mesh as the accuracy control of
self-consistent tunnelling). % <!-- S2PREP:158-160 ; S2PREP:174 -->
A regression script rebuilds every launched deck from its recorded overrides and must report no unexpected
difference before the next launch. % <!-- CD:290 -->

\subsubsection{Physical model}
\label{sec:m-tcad-model}

The campaign model (V3) uses one parameter set for the 2.0, 6.3 and \SI{13.2}{nm} films; film dependence enters only
through the thickness $t$ in physical laws. % <!-- CD:24-27 ; CD:487-489 -->
Its ingredients and their labels are listed in \cref{tab:m-tcad-params}.
\begin{itemize}
\item \emph{Confinement.} The conduction-band shift $\dEc(t)$ (in \texttt{AFFINITY}), the gap opening
$\Delta E_g(t)$ (in \texttt{EG300}) and the mass increment $\Delta m^*(t)$ (in \texttt{NC300}) follow power laws fitted
to our crystalline PBE slab points at H--H thicknesses 0.695, 0.930 and \SI{1.962}{nm}; no device data enter these
laws (\prov{FITTED} to \prov{COMPUTED} points). % <!-- V3:58-74 ; V3:84-97 ; V3:154-172 -->
Values at 6.3 and \SI{13.2}{nm} are extrapolations, and identifying the device thickness with the slab H--H thickness
is \prov{ASSUMED}. % <!-- CD:151 -->
The \SI{0.695}{nm} point comes from a slab whose result is still marked preliminary and is not yet entered in the DFT
results log \pdPending{P02}{CHECK PENDING: enter the 0.75 nm (H--H 0.695 nm) slab in RESULTS\_LOG.md, or refit the
laws without it; claim C021}. % <!-- LEDGER:120-124 ; V3:68 -->
\item \emph{Quasi-two-dimensional density of states.} At \SI{2.0}{nm} only, $\Nc$ is replaced by $N_{2D}/t$
(\prov{CALCULATED}) under a pre-registered single-subband rule (DFT subband spacing \SI{0.7191}{eV} $\gg \kT$). This
substitution inside the 3-D Fermi--Dirac form overstates the degenerate on-state density by up to $1.7\times$. % <!-- V3:98-105 ; CD:63 ; CD:111 -->
\item \emph{Traps.} An exponential acceptor tail with integrated density $N_s$ spread uniformly over the film
($\Nt = N_s/t$, $\NTA = \Nt/\WTA$); a deep acceptor Gaussian; and an interface acceptor sheet at the IWO/\AlO{}
interface. Trap levels are referenced to $\Ec$. % <!-- V3:106-118 ; V3:178-195 -->
\item \emph{Charges and donors.} A uniform front fixed charge $\Qf$ at $y = 0$ carries the alignment offset
$\Delta V_0$; it is degenerate with the gate work function and the bulk affinity, which are therefore held fixed.
The donor density is constant: the fitted V1 amplitude is kept, and the constant form is \prov{ASSUMED}. % <!-- CD:50 ; CD:116 ; V3:119-130 -->
\item \emph{Mobility.} A constant band mobility per film from one shared Matthiessen law,
$\muband(t) = \mu_\infty\,(0.208\,m_0/\mstar(t))/(1+\ell/t)$, multiplied by the roughness factor $(1-\Dsr/t)^2$ with
$\Dsr = \SI{0.287}{nm}$ from the source paper. % <!-- V3:131-153 -->
\end{itemize}
Models switched off, each with a stated reason, are field-dependent mobility, the silicon low-field and inversion-layer
mobility laws, band-to-band tunnelling, Coulombic trap barrier lowering, lumped contact resistance, a back-surface
charge sheet and the density-gradient model. % <!-- CD:129-137 -->
Branches tested once and adopted only by a pre-registered rule are the IGZO.TOKYO carrier-density mobility, Schottky
Pd contacts with a shared barrier from $\{0.2, 0.5\}$~eV, and a Bohm-potential pair test at \SI{2.0}{nm}. % <!-- CD:118-120 ; CD:308-321 -->

\begin{table}[htbp]\centering\footnotesize
\caption{Campaign model V3. Values at 2.0 / 6.3 / \SI{13.2}{nm}. Fitted values are the S4 start values; they are not
calibration results. Labels follow the claims-ledger vocabulary.}
\label{tab:m-tcad-params}
\begin{tabularx}{\linewidth}{@{}>{\raggedright\arraybackslash}p{3.3cm}>{\raggedright\arraybackslash}X>{\raggedright\arraybackslash}p{4.6cm}@{}}
\toprule
Parameter & Value & Label and basis \\
\midrule
Bulk gap, affinity & 3.05~eV; 4.3~eV & \prov{LITERATURE} (another sputtered IWO device) \\ % <!-- V3:27-37 -->
Bulk mass & $0.208\,m_0$ & \prov{LITERATURE} (optical Hall, \InO{}) \\ % <!-- V3:38-42 -->
$\dEc(t)$ & $0.70147\,t^{-1.31669}$~eV: 0.2816 / 0.0622 / 0.0235~eV & \prov{FITTED} to \prov{COMPUTED} DFT points \\ % <!-- V3:155-159 -->
$\Delta E_g(t)$ & $0.81431\,t^{-1.31518}$~eV & \prov{FITTED} to \prov{COMPUTED} DFT points \\ % <!-- V3:60-63 -->
$\mstar(t)$ & $0.208 + 0.11116\,t^{-1.26387}$: 0.2543 / 0.2189 / 0.2123~$m_0$ & \prov{LITERATURE} + \prov{FITTED} increment \\ % <!-- V3:84-89 -->
$\Nc$ (\SI{300}{K}) & $1.373\times10^{19}$ (quasi-2D) at \SI{2.0}{nm}; 3-D formula otherwise & \prov{CALCULATED} \\ % <!-- V3:98-105 -->
Tail $N_s$, $\WTA$ & $4\times10^{12}$\cmsq{}, \SI{40}{meV} (start) & \prov{FITTED} shared (S4) \\ % <!-- V3:106-116 -->
Deep acceptor Gaussian & $5\times10^{16}$~cm$^{-3}$\,eV$^{-1}$ at $\Ec - 0.6$~eV, width 0.15~eV & \prov{ASSUMED}; zero in S6 \\ % <!-- V3:178-182 -->
Interface acceptor sheet & $3\times10^{11}$~cm$^{-2}$\,eV$^{-1}$ at $\Ec - 0.3$~eV, width 0.12~eV & \prov{ASSUMED}; zero in S6 \\ % <!-- V3:189-195 -->
Donors $\Nd$ & $2.5\times10^{17}$\cmcu{}, constant & \prov{FITTED} (V1 amplitude, kept); form \prov{ASSUMED} \\ % <!-- V3:119-128 -->
Front charge $\Qf$ & $1.38345\times10^{12}$\cmsq{} (start) & \prov{FITTED} shared, profiled \\ % <!-- V3:196-199 -->
Mobility law & $\mu_\infty = 100.0$\cmVs{}, $\ell = 7.25$~nm (start) & \prov{FITTED} shared, profiled \\ % <!-- V3:147-153 ; CD:178-179 -->
Capture cross-sections, SRH lifetime & $10^{-15}$~cm$^2$; $10^{-6}$~s & \prov{ASSUMED} \\ % <!-- V3:176-177 ; V3:54-57 -->
Valence-band DOS, hole mobility & $10^{19}$\cmcu{}; 0.01\cmVs{} & \prov{ASSUMED} (holes not solved) \\ % <!-- V3:47-53 -->
\bottomrule
\end{tabularx}
\end{table}

\subsubsection{Model history: per-film calibration and the registered held-out test}
\label{sec:m-tcad-history}

Model V1 was calibrated per film: $\Qf$ to the constant-current threshold and the band mobility to $\Ion$, with the
tail width fitted on the \SI{2.0}{nm} subthreshold region (\prov{FITTED}). % <!-- CS:141-148 -->
Model V2 replaced the literature confinement laws by our DFT laws. Its pass/fail criteria were committed before the
launches (commit \texttt{7b75ee6}); one error in the registered transformation was corrected and published
(\texttt{a914165}) before the held-out outcome existed, and both versions were scored (\texttt{852dc8a}). % <!-- CS:164 ; V2TEX:43-50 -->
The held-out \SI{6.3}{nm} threshold prediction failed in both variants (predicted $\Vthcc$ 0.270 and 0.417~V against
0.644~V measured), as did the earlier shared-parameter V1 test (claim C014, \prov{FALSIFIED}). % <!-- LEDGER:80-84 ; CS:173-178 -->
The V2 off-state strip, a gate-independent parallel conductance with one geometric factor fitted on \SI{13.2}{nm}, is
a \prov{HYPOTHESIS} (claim C006). % <!-- LEDGER:34-38 ; V2TEX:137-145 -->
V1 and V2 enter this paper only as the documented starting point of the campaign below.

\subsubsection{Extraction protocol}
\label{sec:m-tcad-extraction}

Threshold, subthreshold swing and mobility are extracted by one code path for measured and simulated curves, under a
protocol written and committed (\texttt{7a61687}) before any measured-versus-TCAD number was computed. % <!-- PROT:3 ; CS:86 ; DLOG:16 -->
\begin{itemize}
\item \emph{Threshold (primary).} $\Vthcc$ is $\Vg$ at $\Id = \SI{1e-9}{\Aum}$, by log-linear interpolation at the
last upward crossing below +3~V; \SI{5e-10}{\Aum} is also reported. % <!-- PROT:32 -->
\item \emph{Subthreshold swing (registered window).} One decade between $5\times10^{-11}$ and \SI{5e-10}{\Aum},
admissible when the lower edge is at least ten times the floor; otherwise a floor-corrected value is primary, which
assumes an additive, gate-independent parallel current. % <!-- PROT:38-43 -->
\item \emph{Amendment A1 (secondary window).} Because the films differ in current scale, a window normalised by each
film's measured on-current, $[r, 10r]\times\Ion$ with $r = 10\times\max_t(\mathrm{floor}/\Ion) = 2.15\times10^{-5}$,
was added (commit \texttt{7ce7116}) after the measured values under the registered rule had been computed and before
any simulated value under the new rule existed; the disclosure is part of the amendment. % <!-- PROT:80-98 ; A1R:3 ; CS:87 -->
\item \emph{Mobility.} $\muFE = L\,\gm/(W\Cox\Vd)$ from central-difference $\gm$, reported at 1.0--3.0~V; it equals a
carrier mobility only where $\Vg - V_T \ge \Vd$ and the mobility is gate independent, and is otherwise a normalised
transconductance. A gradual-channel power-law fit $\mu = \mu_0 (u/1\,\mathrm{V})^\gamma$ over 1.0--3.0~V gives a
model-based channel mobility. % <!-- PROT:50-58 -->
\item \emph{Uncertainties and comparison.} Measurement noise from local quadratic residuals, interpolation-method
differences and the TCAD numerical budget are combined in quadrature; a difference is ``consistent'' if it is within
$2\sigma$. $\Cox$, the charge-centroid ratio $C_{\mathrm{eff}}/\Cox = 0.87$, contact resistance and device-to-device
spread are stated separately and not included. % <!-- PROT:63-78 -->
\end{itemize}
The source paper's peak mobilities are reproduced numerically on our curves by the peak of the saturation-formula
mobility (claim C030, \prov{DERIVED}); that this is the paper's method is only the leading candidate (C031,
\prov{INFERRED}), and the paper's near-thermal swing at \SI{2}{nm} remains \prov{UNRESOLVED} (C025). We report our
definitions throughout. % <!-- LEDGER:178-191 ; LEDGER:146-152 -->
The \SI{13.2}{nm} swing mismatch is reported as a finding; no window or floor treatment was chosen to remove it
(decision D-06). % <!-- DLOG:17 -->
\pdPending{P03}{CHECK PENDING: the fabrication group's mobility and SS recipe (formula, $\Cox$, $W/L$, current
normalisation, sweep provenance); claims C024, C025, C031}.

\subsubsection{Campaign design}
\label{sec:m-tcad-campaign}

The campaign (C10) was designed before its first launch from three independent inputs (the DFT-to-TCAD extraction, an
audit of the deck generator and ATLAS 5.28, and an independent design by Codex), merged, and reviewed read-only by
Codex. % <!-- CD:3-10 -->
Its budget was 200 sequential launches (decision D-13). % <!-- DLOG:23 ; CD:39-41 -->
\begin{itemize}
\item \emph{Shared fitted parameters.} Five values for all three films: $\Delta V_0$ (written as one $\Qf$), $\mu_\infty$,
$\ell$, $N_s$ and $\WTA$, plus a sixth only if a mobility branch is adopted. No per-film charge or mobility exists. % <!-- CD:24-27 -->
\item \emph{Exact profiling.} A uniform front charge shifts a transfer curve rigidly, and with a constant (or
IGZO-type) mobility the current is proportional to \texttt{MUN}; both hold in the continuum equations up to the
generation term, bounded at \SI{1e-28}{\Aum} against currents of at least \SI{1e-15}{\Aum} (\prov{CALCULATED}). % <!-- CD:228-233 -->
$\Delta V_0$, $\mu_\infty$ and $\ell$ are
therefore profiled offline on each launched film triple, without extrapolation beyond the solved $\Vg$ range, and every
optimum is confirmed natively within \SI{0.01}{dec}. % <!-- CD:28-32 ; CD:225-243 -->
\item \emph{Objective.} $J = \tfrac13\sum_t (\mathrm{RMSE}^2_{\mathrm{sub},t} + \mathrm{RMSE}^2_{\mathrm{on},t})/2$,
with log residuals over the measured subthreshold range from ten times the floor to \SI{1e-9}{\Aum} and over
$\Vg = 1$--3~V; no film is weighted by current or point count. % <!-- CD:215-223 -->
\item \emph{Acceptance} (\prov{ASSUMED} engineering limits, not confidence intervals). Per film: active log-RMSE
$\le \SI{0.15}{dec}$, $|\Delta\Vthcc| \le \SI{0.10}{V}$, $\Id(3\,\mathrm{V})$ within $\pm30\,\%$, SS in the A1 window
within $\pm20\,\%$, $\muFE$ at 1.5, 2 and 3~V within $\pm20\,\%$. Success on all three films is still called
calibration. % <!-- CD:246-255 -->
\item \emph{Identifiability.} A scaled residual Jacobian from the calibration archive; a condition number above 100 or
$|\mathrm{corr}| > 0.95$ (\prov{ASSUMED}) means only a parameter combination is reported. % <!-- CD:257-261 -->
\item \emph{Stages.} S0 smoke (2 launches), S1 baselines and numerics (7), S2 off-floor mechanisms (17), S3 model-form
branches (27), S4 calibration (45 space-filling + 18 adaptive + 3 native confirmation + 3 runner-up), S5
leave-one-film-out (30, retrospective), S6 sensitivities and frozen predictions (27), S7 \SI{31.8}{nm} stress (2),
reserve ($\le 19$). % <!-- CD:269-282 -->
\item \emph{Stop rules.} Among them: if the best $\sqrt J$ after S4a exceeds \SI{0.30}{dec} with systematic residuals,
the shared model is reported as failed and no per-film rescue parameter is introduced; a held-out film that fails the
acceptance falsifies transferability of the model form and is not repaired. % <!-- CD:324-338 -->
\end{itemize}

\paragraph{Off-state floor tests.} Gate leakage is computed in the solver with self-consistent tunnelling
(\texttt{QTNLSC.EL}), a second Fowler--Nordheim formulation (\texttt{FNORD}) and defect-assisted oxide tunnelling
(\texttt{ITAT.SC.EL}); generation is bounded with extreme SRH and trap-assisted tunnelling in two-carrier runs; and an
ungated IWO strip is simulated as a separate deck. % <!-- CD:34-38 ; CD:121-128 -->
Two oxide parameter sets are used: a nominal set from the literature (\HfO{}/\AlO{} affinities 2.5/1.25~eV, gaps
5.8/8.8~eV, tunnelling masses \prov{ASSUMED} inside literature ranges) and a deliberately leakage-maximising set
(TiN/\HfO{} barrier 1.5~eV, \prov{ASSUMED}), which is an optimistic scenario, not a rigorous upper bound. % <!-- CD:200-206 ; CD:82 ; CD:376 ; PHYS:58-78 -->
Because \texttt{FNORD} evaluates the field in the \AlO{} at the semiconductor interface, the \HfO{}-barrier constants
are transferred to that field ($r = \varepsilon_{\mathrm{HfO_2}}/\varepsilon_{\mathrm{Al_2O_3}} = 2.174$). % <!-- CD:81 ; PHYS:34-40 -->
\pdPending{P04}{CITATION PENDING: oxide band parameters (Robertson 2000, J. Vac. Sci. Technol. B 18, 1785, as cited in
the campaign design), tunnelling-mass sources, Lenzlinger--Snow constants, Ielmini trap-assisted tunnelling,
IGZO.TOKYO (Abe 2014); not yet in the bibliography}.
A mechanism explains a film's floor only if at least 80\,\% of the window points are valid, its own contribution
(against a matched control) lies within \SI{0.5}{dec} of the floor, its slope is within twice the measured slope plus
$2\sigma$, and the current balance passes; it explains the thickness scaling only if the three films agree within
\SI{0.5}{dec}. At \SI{2.0}{nm} an under-prediction does not falsify a mechanism, because that floor may be
instrument-limited. % <!-- CD:358-371 -->

\subsubsection{Pre-registered TCAD predictions}
\label{sec:m-tcad-register}

Before any S2 launch, the hypothesis register (\cref{sec:m-reg}) froze which launches score which hypothesis, the
predictions, the stop rules and the calibration membership. % <!-- REG:1089-1095 ; DLOG:28 -->
\Cref{tab:m-tcad-reg} summarises the registered TCAD predictions; every one is scored and reported whatever the
outcome. The diagnostic set comprises 26 launches, with the oxide-mesh check S2-11 first. % <!-- REG:1031-1045 ; REG:1054-1064 -->

\begin{table}[htbp]\centering\footnotesize
\caption{Registered TCAD predictions (register revision 2 and its amendments). Every threshold is an \prov{ASSUMED}
decision threshold. Outcomes are not part of this draft.}
\label{tab:m-tcad-reg}
\begin{tabularx}{\linewidth}{@{}l>{\raggedright\arraybackslash}X>{\raggedright\arraybackslash}p{3.4cm}@{}}
\toprule
ID & Registered prediction & Outcome \\
\midrule
H-B & Strip $I(0.7\,\mathrm{V})/I(0.1\,\mathrm{V}) = 7.0\pm0.5$; with $(W/L)_{\mathrm{par}}$ fitted on \SI{13.2}{nm} only, the floors at 6.3 and \SI{2.0}{nm} within \SI{0.5}{dec} of the measured ones & \pdPending{P05}{RESULT PENDING: S2-14..S2-17, STRIP-6p3} \\ % <!-- REG:126-131 -->
H-C & Nominal leakage at $-3$~V $\le 10^{-2}\times$ floor (only with a demonstrated ceiling below that bound); extreme generation $\le \SI{1e-20}{\Aum}$ or below resolution; worst-case leakage fails flatness or thickness scaling; resolved leakage changes by $\le\SI{0.1}{dec}$ under the S2-11 oxide-mesh refinement & \pdPending{P06}{RESULT PENDING: S2-11 (sealed), S2-01..S2-13} \\ % <!-- REG:139-151 -->
H-D & Band-mobility ratio \SI{2.0}{nm}/\SI{13.2}{nm}: mass-only 0.75--0.95 (reported ``unresolved: not computed'', A22); scattering branch 0.20--0.35 with $\ell = 4$--\SI{12}{nm} (crystalline mass inputs) & \pdPending{P07}{RESULT PENDING: S4, native triples} \\ % <!-- REG:164-176 -->
H-E & All confinement inputs off: $\Delta\Vthcc = -0.13\pm0.03$~V at \SI{2.0}{nm} (A13; the original $-0.23\pm0.04$~V is also scored), $-0.062\pm0.015$~V at 6.3, $-0.023\pm0.008$~V at \SI{13.2}{nm}; S3-01 $-$ S1-05 $= -0.272\pm0.010$~V; affinity-only shifts $-281.6\pm5$ and $-62.2\pm5$~mV, rigid within \SI{5}{mV} & \pdPending{P08}{RESULT PENDING: S3-01..S3-03, AFF-2p0, AFF-6p3} \\ % <!-- REG:184-204 -->
H-F & Two-film reference: \SI{6.3}{nm} needs a rigid shift of $+0.28\pm0.08$~V and a current scale of 0.15--0.35; baseline shape RMSE $> \SI{0.15}{dec}$ & \pdPending{P09}{RESULT PENDING: S3-10..S3-13, S3-27; H-F scoring not yet frozen} \\ % <!-- REG:243-247 ; S2PREP:167-173 -->
H-J & Every film: RMSE $\le\SI{0.15}{dec}$, $\Vth$ within $\pm0.10$~V, $\Ion$ within $\pm30\,\%$, SS and $\muFE$ within $\pm20\,\%$; any failure falsifies transfer & \pdPending{P10}{RESULT PENDING: S4a, S4c, S5} \\ % <!-- REG:286-293 -->
\bottomrule
\end{tabularx}
\end{table}

\subsubsection{Scoring and the resolution method}
\label{sec:m-tcad-scoring}

All registered TCAD scores come from one frozen scorer (\file{score_register.py}, sha256 \texttt{c0f28bb9}\ldots
\texttt{b55a}; self-test 90/90). % <!-- S2FIX:349 ; S2GATE:222 ; S2GATE:243 -->
It applies the register's censoring rule: a quantity below the demonstrated numerical resolution is reported as
\emph{below resolution}, counts as consistent with an upper bound only if that resolution is itself below the bound, and
is otherwise \emph{unresolved}, never a pass. % <!-- REG:16-20 -->
The register does not define ``demonstrated resolution''. The scorer therefore uses an operational method, which we
label \prov{ASSUMED}:
\begin{itemize}
\item a resolution is demonstrated only by a pair of two distinct completed runs that differ by one verified numerical
intervention (trap-DOS levels, global mesh factor, or oxide $y$-mesh factor), checked against the actual deck values
and recomputed from the run files, which are bound by sha256, at every use; % <!-- S2FIX:20 ; S2FIX:150 -->
\item the pair must agree at least at three points above the candidate resolution, one of them within one decade, over
the full bias range, with a deck-ratio tolerance of $10^{-3}$; % <!-- S2FIX:162-170 -->
\item a resolution applies only to its own quantity, film, physical settings, bias range and numerical endpoint; it is
never transferred between films, and for an upper-bound test the pair may not contain a tested run; % <!-- S2FIX:156-160 ; S2FIX:273 -->
\item a verdict computed with the runner's current-balance limit as a proxy is labelled \emph{provisional} and can
never pass an upper bound. % <!-- S2PREP:122 ; S2FIX:23 -->
\end{itemize}
Pair agreement establishes numerical stability under the tested intervention, not an absolute error bound. The
one-decade support value was chosen knowing the two numerical pairs already seen in S1; it is prospective for S2 only. % <!-- S2FIX:288-291 ; S2FIX:167 -->
With the current run plan, almost every primary S2 and affinity-only verdict will therefore be provisional or
unresolved. % <!-- S2FIX:238-247 -->
\pdPending{P11}{DECISION PENDING: the user's sign-off of the resolution method as a logged amendment, required before
the S2-11 outputs are unsealed}. % <!-- S2FIX:129 ; S2GATE:248 -->
\pdPending{P12}{DECISION PENDING: whether to add matched numerical refinement runs (a logged plan change) so that
primary S2 verdicts can become demonstrated}. % <!-- S2FIX:247 -->
The \SI{6.3}{nm} shape-residual hypothesis (H-F) has no frozen operational scoring yet; its five launches are held. % <!-- S2PREP:137 ; S2PREP:167-173 -->
The oxide trap-assisted tunnelling runs (S2-12, S2-13) cannot count as converged without a further logged amendment
(for example a \texttt{TAT.SLICEPTS} sensitivity launch), because the S2-11 verdict covers \texttt{QTNLSC} only. % <!-- S2PREP:164 -->

\subsubsection{Execution controls, sealing and protocol deviations}
\label{sec:m-tcad-controls}

The campaign runner enforces the budget cap, sequential execution and refusal of held plans in code. % <!-- CD:288-289 ; S2PREP:37 -->
For the diagnostic stage, executable stop gates evaluate the frozen leakage and generation scorers in an isolated
process and expose only a state (fired, not fired, unresolved); every gate decision is bound by sha256 to its
declaration, the scorer, the scored run files and the cited evidence, and an unresolved or fired gate stops the plan
until an audit file citing that decision is supplied. % <!-- S2GATE:32-60 ; S2GATE:186-207 -->
A state-only mode keeps metrics, error texts and launch output in a sealed side record, so that run status can be
followed without exposing results; timeout clean-up terminates only identified processes of the launch and never the
licence daemons. % <!-- S2GATE:112-123 ; S2GATE:82-88 ; S2GATE:186-191 -->
\pdPending{P13}{CHECK PENDING: the runner revision 3 (stop gates, sealing, clean-up) was uncommitted and DIAG not
launched as of S2\_RUNNER\_GATES\_2026-10-13.md; record the commit and the Codex verdict}. % <!-- S2GATE:3 -->

The following deviations are disclosed.
\begin{itemize}
\item The smoke run S0-02 needed six launches against a cap of three repairs. % <!-- REG:39 -->
\item The regression gate read five unexpected differences while runs \run{0067} to \run{0074} were launched, against
the rule that it must read zero; the five were reclassified as known after line-by-line verification on 2026-10-11
(review item R23). % <!-- S2PREP:82 -->
\item The registered S1-05 expectation ($+40$ to $+70$~mV) failed ($+142$~mV) and stays recorded as a failure; the
\SI{2.0}{nm} confinement band was re-registered after S1 had been seen and before any S3 run, and both bands are scored
(A13). % <!-- REG:31 ; REG:184-190 ; REG:211-212 -->
\item The two-carrier check S1-06 timed out and is unresolved. % <!-- REG:32 -->
\item S2-11 (\run{0075}) has been launched; its outputs are sealed and were not opened for this draft. Its eligibility
under the timeout rule and its ``possibly exposed'' label await rulings. % <!-- S2GATE:11-16 ; S2GATE:123 ; S2GATE:212 ; S2GATE:249 -->
\pdPending{P14}{DECISION PENDING: run\_0075 rulings (timeout eligibility, ``possibly exposed'' label) and the
unsealing of S2-11}.
\end{itemize}

\subsubsection{Launch accounting}
\label{sec:m-tcad-launches}

Every launch counts, failures included. % <!-- CD:287 -->
\Cref{tab:m-tcad-launches} gives the accounting as recorded.

\begin{table}[htbp]\centering\footnotesize
\caption{ATLAS launch accounting (global counter and the C10 campaign counter).}
\label{tab:m-tcad-launches}
\begin{tabularx}{\linewidth}{@{}>{\raggedright\arraybackslash}p{4.6cm}>{\raggedright\arraybackslash}X@{}}
\toprule
Stage & Launches \\
\midrule
V1 report (2026-09-25) & 52 recorded launches \\ % <!-- MAIN:13 ; MAIN:35 -->
Before the campaign (2026-10-07) & 60 of 71 approved \\ % <!-- CS:1 ; CS:276 -->
Campaign approval (D-13) & 200 more; global cap 260 $= 60 + 200$, enforced in code \\ % <!-- DLOG:23 ; SOLV:18 -->
At the register freeze & 14 C10 launches used (\run{0061}--\run{0074}), 186 remaining \\ % <!-- REG:38 -->
As of 2026-10-13 & 75 of 260 global; 15 of 200 C10 (including S2-11, \run{0075}) \\ % <!-- S2GATE:1 ; S2GATE:9 -->
Planned diagnostic set & 26 launches; C10 use 40 of 200 at its end \\ % <!-- REG:1038 ; S2PREP:141 -->
Projected total & about 110--140 launches \\ % <!-- REG:1066 -->
Final count & \pdPending{P15}{RESULT PENDING: final launch count and status of every launch} \\
\bottomrule
\end{tabularx}
\end{table}
